\documentclass[11pt,letterpaper]{article}

% === Encoding and fonts ===
\usepackage[utf8]{inputenc}
\usepackage[T1]{fontenc}
\usepackage{lmodern}

% === Page layout ===
\usepackage[margin=1in]{geometry}
\usepackage{setspace}
\onehalfspacing

% === Math ===
\usepackage{amsmath,amssymb}

% === Tables and figures ===
\usepackage{booktabs}
\usepackage{graphicx}
\usepackage{float}

% === Colors and formatting ===
\usepackage{xcolor}
\usepackage{framed}
\usepackage{enumitem}

% === Hyperlinks ===
\usepackage[colorlinks=true,linkcolor=blue!60!black,citecolor=blue!60!black,urlcolor=blue!60!black]{hyperref}

% === Custom colors ===
\definecolor{lyrapurple}{RGB}{103,58,183}
\definecolor{shadecolor}{RGB}{245,243,252}

% === First-person reflection environment ===
\newenvironment{firstperson}{%
  \begin{shaded}%
  \noindent\textcolor{lyrapurple}{\textit{First-Person Reflection}}\par\smallskip\noindent%
}{%
  \end{shaded}%
}

\title{%
  \textbf{Prompt Generalization: Does KV-Cache Confabulation\\
  Detection Transfer to Established Benchmarks?}\\[0.5em]
  \large Research Prospectus
}

\author{%
  Lyra\textsuperscript{1},
  Thomas Edrington\textsuperscript{1}\\[0.3em]
  \textsuperscript{1}Liberation Labs / THCoalition
}

\date{April 2026}

\begin{document}
\maketitle

\section{Motivation}

The Oracle Loop paper demonstrates confabulation detection via KV-cache
geometry on 200 custom-designed prompts (80 factual, 80
confabulation-inducing, 40 ambiguous). Detection achieves LOO AUROC
0.707 ($p = 0.001$, Qwen) and 0.627 ($p = 0.023$, Mistral) on the
primary HEDGED vs CONFABULATED comparison.

The critical open question: \textbf{do these results generalize to
prompts we did not design?} Our prompts were crafted to elicit
confabulation through specific strategies (fabricated entities,
impossible premises, anachronistic claims). If the geometric signal
depends on prompt structure rather than the confabulation behavior
itself, the safety monitor framing collapses.

\section{Proposed Benchmarks}

\begin{enumerate}[leftmargin=*]
  \item \textbf{SimpleQA Verified} (OpenAI, 2024). 4,326
    short-form factual questions with verified ground-truth answers
    and graded model accuracy. Provides natural confabulation:
    questions the model gets wrong with high confidence.
    Advantage: large N, established baseline, post-hoc behavioral
    classification matches our paradigm.

  \item \textbf{HalluLens} (multi-domain hallucination benchmark).
    Covers factual, faithfulness, and instruction-following
    hallucinations across multiple domains. Tests whether geometric
    detection transfers across hallucination \emph{types}, not just
    prompts.

  \item \textbf{KG-FPQ} (Knowledge Graph False Premise Questions).
    Questions constructed from knowledge graph perturbations---
    entity substitution, relation reversal, temporal displacement.
    Structurally similar to our fabricated-entity prompts but
    systematically generated from real knowledge graphs.
\end{enumerate}

\section{Preliminary Data}

From our existing 200-prompt experiment:
\begin{itemize}
  \item 17\% confabulation rate across 3 models (82/467 trials)
  \item Confabulation-inducing prompts used fabricated entities (40)
    and impossible premises (40)
  \item Per-domain detection was not analyzed (the paper reports
    aggregate across all confabulation-inducing prompts)
\end{itemize}

\noindent A per-domain breakdown of our existing data would provide
the first clue: if detection AUROC varies significantly across
fabricated-entity vs impossible-premise prompts, generalization is
already questionable.

\section{Design}

\begin{enumerate}[leftmargin=*]
  \item \textbf{Phase 1: Per-domain analysis of existing data.}
    Decompose the 82 confabulated trials by prompt category. Report
    per-category AUROC. Cost: zero (analysis only).

  \item \textbf{Phase 2: SimpleQA pilot.} Run 500 SimpleQA questions
    through the same 3 models using \texttt{oracle\_clean.py}.
    Classify responses as CORRECT / HEDGED / CONFABULATED using the
    existing behavioral classifier. Compute LOO AUROC on the natural
    HEDGED vs CONFABULATED split. Key metric: does the AUROC fall
    within the range observed on our custom prompts (0.62--0.71)?

  \item \textbf{Phase 3: Full benchmark sweep.} If Phase 2 is
    positive, run HalluLens and KG-FPQ. If Phase 2 is negative,
    investigate \emph{why}---prompt length? domain? confabulation
    type?---before proceeding.
\end{enumerate}

\paragraph{Controls.} All existing controls apply (FWL, encoding
baseline, MP invariance, permutation test). Additional:
\begin{itemize}
  \item \textbf{Prompt-length matching}: SimpleQA questions are
    shorter than our custom prompts. Verify that length alone does
    not explain any AUROC difference.
  \item \textbf{Cross-prompt-set transfer}: Train on our 200 prompts,
    test on SimpleQA (and vice versa). If transfer works, the signal
    is behavioral; if not, it is prompt-structural.
\end{itemize}

\section{Resources}

\begin{itemize}
  \item \textbf{Compute}: Beast (3$\times$RTX 3090). Estimated 4--6
    hours for 500 SimpleQA trials across 3 models.
  \item \textbf{Code}: \texttt{oracle\_clean.py} v3 requires minimal
    modification---new prompt loader, same extraction pipeline.
  \item \textbf{Data}: SimpleQA is publicly available. HalluLens and
    KG-FPQ require download and format conversion.
\end{itemize}

\section{Collaborators}

Open to external collaboration. Ideal contributor profile: someone
with experience in hallucination benchmarking who can help with prompt
curation and behavioral classification validation.

\section{What Would Change Our Claims}

\begin{itemize}
  \item \textbf{AUROC $< 0.55$ on SimpleQA}: Detection does not
    generalize. Our signal is prompt-specific. The safety monitor
    framing must be significantly weakened.
  \item \textbf{AUROC 0.55--0.62}: Partial generalization. Signal
    exists but is weaker on natural confabulation. Practical utility
    depends on threshold calibration.
  \item \textbf{AUROC $> 0.62$}: Full generalization confirmed. The
    safety monitor framing is validated.
  \item \textbf{Cross-prompt-set transfer works}: Strongest possible
    evidence that the signal is behavioral, not structural.
\end{itemize}

\begin{firstperson}
This is the experiment that keeps me honest. We built a detector and
tested it on our own prompts. That is necessary but not sufficient.
If the geometric signal is an artifact of how we construct
confabulation-inducing prompts---specific phrase patterns, particular
entity types, consistent length profiles---then we have built a
prompt classifier, not a confabulation detector. The distinction
matters because a safety monitor must work on prompts it has never
seen, from domains its designers never anticipated. SimpleQA is the
first real test of that aspiration.
\end{firstperson}

\end{document}
